\section{A fixed-energy fractional-moment estimate}
\label{sec:bootstrap}

We return to the true product law and remove the hierarchy from the
operator whose decay is to be estimated. Throughout this section,
\(h>0\) and \(E\in(-4-h,4+h)\) are fixed. All choices may depend on
these two parameters. Distance means the shortest-coordinate max-norm
distance on the torus \(\T_N^2=(\Z/b_N\Z)^2\). The notation for the probe
sets below is local to this section.

\begin{theorem}[True torus Green function]\label{boot:fractional}
Fix \(0<s<1/16\). There are constants
\(A=A(h,E,s)<\infty\), \(a=a(h,E,s)>0\), and an integer
\(N_0=N_0(h,E,s)\) such that, for every \(N\ge N_0\) and every
\(x,y\in\T_N^2\),
\begin{equation}\label{boot:real-decay}
 \E\abs{(H_N-E)^{-1}_{xy}}^{s/2}
       \le A\exp\{-a\dist(x,y)\}.
\end{equation}
Here \(H_N\) has the true iid uniform \([-h,h]\) diagonal and periodic
nearest-neighbor hopping. Its inverse at the specified real energy
exists almost surely. No uniformity in \(E\), nor any statement about
time evolution, is part of this theorem.
\end{theorem}

We prove the theorem by first retaining arbitrary deterministic masks of
artificial probes. At each site a mask prescribes absence, strength
\(\tau\), or strength \(\tau\vartheta(v_i)\), with at most one probe per
site. The profile \(\vartheta\) is the one fixed earlier; in particular
it equals one on \([-h/2,h/2]\), is \(C^{1,1}\), and vanishes quadratically
only at the endpoints of the potential interval. Write
\(\eta_i\) for a prescribed positive probe strength and let \(Y\) have
columns \(\sqrt{\eta_i}e_i\). In this section
\[
 G=(H_N-E-iYY^\top)^{-1},\qquad S=1+2iY^\top GY.
\]
Thus \(S\) is symmetric and unitary, and \(Y^\top GY\) is a
contraction by \Cref{mat:ward}. All exceptional algebraic singularities
below can be discarded under the continuous product law.

The parameter order is
\begin{equation}\label{boot:order}
 s,\ d_0,\ C_0,\ n,\ m,\ N,
 \qquad R=\lfloor\exp(b_m^2)\rfloor,\qquad \tau=C_0/R.
\end{equation}
The barrier threshold \(d_0>0\) will be small; \(C_0\) will be large.
After they have been fixed, \(n\) and then \(m>n\) are made sufficiently
large. Finally \(N\) is sufficiently large in terms of all these
choices; all ensuing estimates hold for every larger \(N\). We always
require \(b_N>100R\) and \(\tau\le1\).

\subsection{Filling and twisting an annulus}

Let \(X,Y'\) be arbitrary deterministic subsets of the probe sites in
the radius-\(R\) squares centered at \(x,y\), respectively, where
\(\dist(x,y)\ge10R\). Let
\[
 C=\{i:2R\le\dist(i,x)\le4R\}.
\]
Fill every previously unoccupied site of \(C\) with a probe of strength
\(\tau\vartheta(v_i)\), keeping existing probes and their strengths.
The filled mask is again deterministic. Superscripts \({\rm old}\) and
\({\rm fill}\) distinguish the resulting scattering matrices. Identify
the probes in \(C\) with their sites and put
\[
 U=\norm{S^{\rm fill}_{CX}}_{\rm HS},\qquad
 V=\norm{S^{\rm fill}_{CY'}}_{\rm HS},\qquad
 D=\norm{S^{\rm fill}_{Y'X}}_{\rm HS}.
\]

\begin{lemma}[Mixture domination]\label{boot:twist}
There is \(C_s<\infty\), independent of the masks, centers, \(R,N\),
and \(C_0\ge1\), such that
\begin{equation}\label{boot:twist-ineq}
 \E\norm{S^{\rm old}_{Y'X}}_{\rm HS}^{s}
     \le C_s(1+C_0)\E\bigl[D^s+V^sU^s\bigr].
\end{equation}
As everywhere in this section, the constant may depend on the fixed
\(h,E,s\) and profile.
\end{lemma}

\begin{proof}
Close just the newly added probes using a common Haar phase \(z\).
The closure identity in \Cref{mat:closure} and the sandwich estimate in
\Cref{mat:weak}, followed by integration of the weak \(L^1\) tail, give
\[
 \int\norm{(S_z)_{Y'X}}_{\rm HS}^{s}\,\dd z
       \le C_s(D^s+V^sU^s).
\]
Here \(\dd z\) is normalized Haar measure, and the blocks linking the
new probes to \(X,Y'\) are subblocks of those defining \(U,V\).
The corresponding real closing shift at each newly occupied site is
\(w\vartheta(v_i)\), where \(w\) has density
\[
                 \frac{\tau}{\pi(w^2+\tau^2)}.
\]
Only newly occupied sites are shifted, so after closing we have precisely
the old mask evaluated at these shifted true potentials.

For \(\abs w\le c/R\), with fixed sufficiently small \(c>0\), the map
\(T_w(v)=v+w\vartheta(v)\) is a diffeomorphism of \([-h,h]\) onto
itself. If \(u=T_w(v)\), its pushforward density relative to the uniform
law is \(r_w(u)=(1+w\vartheta'(v))^{-1}\). Uniformly in \(u\),
\[
 \log r_w(u)=-w\vartheta'(u)+O(w^2).
\]
This follows from the inverse Jacobian and the Lipschitz continuity of
\(\vartheta'\). There are at most \(K R^2\) shifted sites, so at each
fixed output configuration the joint density satisfies
\[
 \log\prod_i r_w(u_i)
        =-w\sum_i\vartheta'(u_i)+O(R^2w^2).
\]
Choose the half interval \([0,c/R]\) or \([-c/R,0]\) on which the linear
term is nonnegative. On this entire half interval the joint density is
at least a fixed positive constant. Its Cauchy probability is
\[
 \frac1\pi\arctan\frac c{C_0}\ge\frac{c_1}{1+C_0}.
\]
Consequently the full shifted mixture dominates
\(c_2(1+C_0)^{-1}\) times the true product law, pointwise in the output
configuration. Applying this domination to the nonnegative old-mask
block power and then averaging the Haar estimate proves the claim.
\end{proof}

We seek a bound by an arbitrarily small multiple of \(\E V^s\).
Thus \Cref{boot:twist-ineq} leaves two tasks: attenuate \(D\) across
the filled annulus, and combine a small mean return \(U\) with a
resampling bound for \(U^sV^s\). An unweighted estimate for barrier
failure would not suffice when \(\E V^s\) is small. We therefore
control that failure also under the \(V^s\)-weighted law.

\subsection{Two barriers under a transmission tilt}

Call a site good when \(\vartheta(v_i)\ge d_0\). Give cost one to an
edge precisely when both endpoints are good and both radii about \(x\)
belong to \([2.2R,2.8R]\); all other edges have cost zero. Let
\(d_{\rm in}\) be the associated path-cost distance from the radius-\(2R\)
square. Independently of this definition, count good-good edges with
both radii in \([3.2R,3.8R]\), and let \(d_{\rm out}\) be the resulting
distance from \(\{i:\dist(i,x)\ge4R\}\). Integer roundings of the four
band boundaries make no difference to the constants. For a sufficiently
small absolute \(\gamma>0\), define
\begin{equation}\label{boot:barrier-event}
 \mathcal B=\left\{
 \begin{array}{ll}
 d_{\rm in}(i)\ge\gamma R&\text{if }\dist(i,x)\ge3R,\\
 d_{\rm out}(i)\ge\gamma R&\text{if }\dist(i,x)\le3R.
 \end{array}\right\}.
\end{equation}
These two strictly buffered bands will allow separate resampling of the
inside and outside of radius \(3R\).

\begin{lemma}[Uniform one-site estimates under tilt]\label{boot:tilted}
Suppose \(\E V^s>0\), and define the probability law
\[
                      \dd\nu=\frac{V^s}{\E V^s}\,\dd\PP.
\]
For a site \(i\in C\), conditional on all other potentials, put
\(W(v_i)=V^s\). Then
\begin{equation}\label{boot:l2l1}
       \bigl(\E_i W^2\bigr)^{1/2}\le K_s\E_i W,
\end{equation}
with a constant uniform in every frozen exterior and in the mask.
For \(d_0>0\) sufficiently small, chosen after \(s\) and before \(C_0\),
there are \(K,c>0\) such that
\begin{equation}\label{boot:tilted-barriers}
 \PP(\mathcal B^c)+\nu(\mathcal B^c)\le K e^{-cR},
 \qquad
 \E_\nu d_{\min}^{-2s}\le K R^2,
 \qquad d_{\min}=\min_{i\in C}\vartheta(v_i).
\end{equation}
No independence under \(\nu\) is asserted or needed.
\end{lemma}

\begin{proof}
Fix all variables except \(v=v_i\), and write the probe factor at \(i\)
as \(a(v)\in\{1,\vartheta(v)\}\). Schur complementation at \(i\)
expresses the other rows and the row at \(i\) of the block defining \(V\)
in the form
\[
 A'+\frac{D'}{v-i\tau a(v)-w_*},\qquad
 \frac{D''\sqrt{a(v)}}{v-i\tau a(v)-w_*},\qquad \Im w_*\ge0.
\]
The coefficients are fixed matrices of compatible sizes. Passivity
gives the sign of \(\Im w_*\). The deleted physical-site system is
invertible almost surely: the imaginary part forces a kernel vector to
vanish at all remaining probe sites, so it would lie in the kernel of
the true real restriction, whose nonzero determinant polynomial
vanishes only on a product-law null set.
Writing both rows over their common denominator, their joint numerator is
\[
 N(v)=\bigl((v-i\tau a(v)-w_*)A'+D',\ D''\sqrt{a(v)}\bigr).
\]
On the plateau \(P=[-h/2,h/2]\), this is the affine vector
\(N_P(v)=(vA'+B,D'')\), where \(B=D'-(i\tau+w_*)A'\).
If \(M=\max_{v\in P}\norm{N_P(v)}_{\rm HS}\), then
\[
 \norm{A'}_{\rm HS}\le2M/h,\qquad
 \norm{(B,D'')}_{\rm HS}\le M,\qquad
 \sup_{[-h,h]}\norm{N(v)}_{\rm HS}\le K_hM.
\]
Moreover an interval in \(P\) of length at least \(h/4\), adjacent to
an endpoint attaining the maximum, has numerator norm at least \(M/2\).
This uses the displayed slope bound. If \(M=0\), the entire block
vanishes and the claimed inequality is immediate.

Put \(b(v)=v-i\tau a(v)-w_*\). Its modulus is at most
\(K_h(1+\abs{w_*})\). For every \(0<q<1\),
\begin{equation}\label{boot:denominator}
 \E_i\abs{b(v)}^{-q}\le K_{h,q}(1+\abs{w_*})^{-q}.
\end{equation}
Indeed, for bounded \(w_*\) use
\(\abs{b(v)}\ge\abs{v-\Re w_*}\) and integrate the one-dimensional
singularity; for \(\abs{w_*}\) sufficiently large use
\(\abs{b(v)}\ge\abs{w_*}-h-1\). Hence
\[
 \E_i W\ge c_{h,s}\left(\frac M{1+\abs{w_*}}\right)^s,
 \qquad
 \E_i W^2\le K_{h,s}\left(\frac M{1+\abs{w_*}}\right)^{2s}.
\]
This proves \Cref{boot:l2l1}, uniformly even for arbitrarily large
exterior impedance.

For any measurable single-site set \(A\), Cauchy--Schwarz now gives
\[
 \nu(v_i\in A\mid(v_j)_{j\ne i})
       \le K_s\PP(v_i\in A)^{1/2}
\]
whenever the conditional normalizing mean is positive. The remaining
conditional configurations have zero \(\nu\)-measure. In particular
the right side for \(A=\{\vartheta<d_0\}\) tends to zero with \(d_0\),
uniformly in the conditioning. Iterating this estimate over any fixed
set of distinct sites bounds the probability that all of them are bad
by \(\beta^{k}\), where \(\beta\) can be made arbitrarily small.

A failure of either barrier gives, after trimming to intermediate
radii and erasing loops, a simple path in the corresponding band of
length \(j\ge c_3R\) and cost less than \(\gamma R\). For \(\gamma\)
small, a fixed positive fraction of its edges are not good-good.
Since the path is simple, this forces at least \(c_4j\) distinct bad
sites. There are at most \(K R^2 4^j\) candidate paths of length \(j\),
and at most \(2^{j+1}\) subsets of their vertices. Choose \(d_0\) so
that \(8\beta^{c_4}<1\). Summing the resulting geometric series for
\(j\ge c_3R\) proves the tilted exponential bound. The same argument
with the original independent site probabilities proves the untilted
bound.

Finally, the same conditional Cauchy--Schwarz estimate gives
\[
 \E_\nu[\vartheta(v_i)^{-2s}\mid(v_j)_{j\ne i}]
       \le K_s\bigl(\E_i\vartheta(v_i)^{-4s}\bigr)^{1/2}\le K.
\]
The integral is finite because \(\vartheta\) vanishes quadratically and
\(8s<1\). Since
\(d_{\min}^{-2s}\le\sum_{i\in C}\vartheta(v_i)^{-2s}\) and
\(\abs C\le K R^2\), this proves the last assertion.
\end{proof}

\subsection{Cutoffs and the small return estimate}

We record explicitly the consequence of \Cref{van:barrier} that will be
used several times. Suppose \(\Psi=GY_J\) is a matrix of fields sourced
at a probe set \(J\), and a scalar cutoff \(\chi\) changes only across
edges whose two endpoints belong to \(C\) and carry probes of strengths
at least \(\tau d\). If \(\abs{\chi_i-\chi_j}\le L\), then the hopping
commutator can be written as a probe source:
\begin{equation}\label{boot:commutator}
 [H_N,\chi]\Psi=Y_C B_\chi,
 \qquad
 \norm{B_\chi}_{\rm HS}
       \le\frac{K L}{\tau d}
           \norm{Y_C^\top\Psi}_{\rm HS}.
\end{equation}
Only endpoints of changing edges need be included on the right.
Indeed, in the coefficient at an endpoint one divides by its square
root strength, and replaces each neighboring field by its absorbed
field divided by that neighbor's square root strength. Bounded lattice
degree gives the stated operator bound. Ward then yields
\begin{equation}\label{boot:cutoff-ward}
 \norm{Y^\top G[H_N,\chi]\Psi}_{\rm HS}
             \le\norm{B_\chi}_{\rm HS}.
\end{equation}
The output here includes \emph{all} probes. These estimates also compare
two systems agreeing on the cutoff support and its changing edges:
cut off a field of the first system and correct its commutator residual
in the second. No estimate on the exterior resolvent is involved.

\begin{lemma}[Small return to a filled annulus]\label{boot:return}
For each fixed sufficiently large \(C_0\) and each \(\epsilon>0\),
one can choose \(n\), then \(m>n\), then a threshold for \(N\), so that
\[
                          \E U^2<\epsilon
\]
for every subsequent \(N\), every deterministic mask, every center,
and every permitted subset \(X\).
\end{lemma}

\begin{proof}
Carry the filled artificial mask through the auxiliary hierarchy.
Take a scalar voltage \(f_R\) equal to one on the radius-\(R\) square,
zero at radii at least \(2R\), and Lipschitz with constant \(K/R\).
By \Cref{screen:profiles}, at fixed \(n\),
\[
 \E\mathcal Q_{n,N}(f_R)\le Kd_n+o_{m,N;n,C_0}(1).
\]
Apply \Cref{hier:comparison} to thin to uniform stage \(m\).
The profile is supported in the radius-\(2R\) square and has Lipschitz
constant \(K/R\), while \(b_m\ll R\). Thus
\Cref{hier:charges}(ii) bounds its unnormalized charge by \(K r_n\),
where \(r_n=\sum_{l\ge n}q_lp_l\to0\), independently of \(R\).
Writing \(U_m\) for the same artificial-probe block at stage \(m\),
we obtain
\begin{equation}\label{boot:return-stage}
 \E U_m^2
    \le Kd_n+o_{n\to\infty}(1)+o_{m,N;n,C_0}(1).
\end{equation}
Each probe pair contributing to \(U_m^2\) has unit squared voltage difference and is
counted in both directions in \(\mathcal Q\), canceling its factor
\(1/2\). The artificial probes are allowed in the comparison because
\(C_0/R\le\exp(-c_*b_m)\) for sufficiently large \(m\).

Let \(\mathcal A\) be the event that no selected tile of level at least
\(m\) meets the radius-\(5R\) square. Tile counting gives
\[
 \PP(\mathcal A^c)
       \le K\sum_{l\ge m}(1+R/b_l)^2 e^{-b_l^6}
       =o(R^{-4}).
\]
For example each fixed polynomial loss in \(R\) is dominated by
\(\exp(-b_m^6)\), since \(\log R\sim b_m^2\).
On \(\mathcal A\), the stage-\(m\) and true filled coefficients and
probes coincide throughout this square. On \(\mathcal B\) use
\[
                 \chi=(1-d_{\rm in}/(\gamma R))_+.
\]
It equals one on \(X\), vanishes outside radius \(3R\), and changes
only on shared good-good edges. Cut off stage-\(m\) fields sourced at
\(X\), and correct in the true system using
\Cref{boot:commutator,boot:cutoff-ward}. Since \(R\tau=C_0\),
\[
                   U\le\left(1+\frac K{C_0d_0}\right)U_m
                   \quad\hbox{on }\mathcal A\cap\mathcal B.
\]
Always \(U^2\le\abs X\le K R^2\), by unitarity. The exceptional
contribution is therefore \(o(1)\), by the preceding hazard estimate
and the product-law path count in the proof of \Cref{boot:tilted}.
That count does not require \(\E V^s>0\). Now \(d_n\to0\) by
\Cref{van:zero}; choosing \(n,m,N\) in order in
\Cref{boot:return-stage} proves the assertion. The error notation means
that at fixed \(n,C_0\), any desired tolerance holds for all sufficiently
large \(m\), and then for all \(N\) above a threshold depending on \(m\).
It asserts no rate uniform in \(n\).
\end{proof}

\subsection{Direct transmission and independent resampling}

We next prove two pointwise estimates:
\begin{equation}\label{boot:direct}
 D\le Kd_{\min}^{-1}V\quad\hbox{almost surely},\qquad
 D\le K e^{-c C_0}V\quad\hbox{on }\mathcal B.
\end{equation}
For the first, use fields sourced at \(Y'\), symmetry of \(S\), and a
cutoff equal to one inside radius \(2R\), zero outside radius \(4R\),
with gradient at most \(K/R\). Its source term vanishes, while both
endpoints of every changing edge have strength at least
\(\tau d_{\min}\). \Cref{boot:commutator,boot:cutoff-ward} give
\(D\le K(R\tau d_{\min})^{-1}V\); here \(C_0\ge1\).
In particular, if \(\E V^s=0\), then \(D=0\) almost surely, and no
tilted law is required.

For the improved estimate choose an integer
\[
                 \ell=\left\lceil\frac{K_1}{\tau d_0}\right\rceil
\]
with \(K_1\) sufficiently large. For \(j\ge1\) use the cutoff
\[
 \chi_j(i)=\min\left\{1,
     \frac{(d_{\rm out}(i)-(j-1)\ell)_+}{\ell}\right\}.
\]
All original sources have \(d_{\rm out}=0\). Nonzero gradients have
both endpoint distances at least \((j-1)\ell\), because cost distances
are integer valued and change by at most one across an edge.
Gradients only occur on counted good-good edges in the outer band.
For the field \(\Psi=G^{\rm fill}Y_{Y'}\), put
\[
 A_j=\norm{\one_{\{d_{\rm out}\ge j\ell\}}Y^\top\Psi}_{\rm HS},
 \qquad
 B_j=\norm{\one_{C\cap\{d_{\rm out}\ge j\ell\}}
                        Y^\top\Psi}_{\rm HS}.
\]
The cutoff equations and Ward give
\[
             A_j\le\frac K{\ell\tau d_0}B_{j-1}
                    \le\tfrac12 B_{j-1},\qquad B_j\le A_j.
\]
Start with \(B_0=V/2\). On \(\mathcal B\), the set \(X\) has
\(d_{\rm out}\ge\gamma R\), so it remains beyond every one of the
\(\lfloor\gamma R/\ell\rfloor\) cutoffs. For sufficiently large fixed
\(C_0\), followed by sufficiently large \(R\), this number is at least
\(c(d_0)C_0\). The iteration proves the second estimate in
\Cref{boot:direct}. Notice that \(A_j\) includes every absorbed output,
including \(X\); the proof does not require \(X\subset C\).

To separate the product \(V^sU^s\), enlarge the probability space with
two independent fresh families of true potentials. Let
\(A=\{i:\dist(i,x)<3R\}\). Form one comparison configuration by
retaining the original potentials on \(A^c\) and replacing those on
\(A\) by the first fresh family. Form a second by retaining the
original potentials on \(A\) and replacing those on \(A^c\) by the
second fresh family. Keep the filled deterministic mask in both
configurations, and recompute its strengths from the corresponding
potentials. Denote the first configuration's transmission to \(Y'\)
by \(V^{\rm ext}\), and the second's transmission from \(X\) by
\(U^{\rm int}\).

On \(\mathcal B\), the cutoff
\((1-d_{\rm out}/(\gamma R))_+\) is supported in the shared exterior
\(A^c\), equals one on \(Y'\), and has all changing edges in its
buffered outer band. Applying it to the fresh-interior fields and
correcting in the original system proves
\[
                  V\le\left(1+\frac K{C_0d_0}\right)V^{\rm ext}.
\]
Likewise \((1-d_{\rm in}/(\gamma R))_+\) is supported in the shared
interior \(A\) and gives
\[
                  U\le\left(1+\frac K{C_0d_0}\right)U^{\rm int}.
\]
Take \(C_0\) large enough that both factors are at most two. The two
comparison configurations are independent \emph{unconditionally}:
they use disjoint parts of the original product sample and independent
fresh families. Each has the original true product distribution.
Dropping \(\one_{\mathcal B}\) before factoring therefore yields
\begin{equation}\label{boot:factorization}
 \E[\one_{\mathcal B}V^sU^s]
       \le 2^{2s}\E V^s\,\E U^s
       \le 2^{2s}\E V^s(\E U^2)^{s/2}.
\end{equation}
This factorization makes no assertion of independence conditional on
the original barrier event.

The exceptional event is controlled under the tilted law from
\Cref{boot:tilted}. Unitarity gives \(U\le K R\), and
\Cref{boot:direct,boot:tilted-barriers} give, respectively,
\begin{align*}
 \frac{\E[\one_{\mathcal B^c}V^sU^s]}{\E V^s}
       &\le K R^s\nu(\mathcal B^c),\\
 \frac{\E[\one_{\mathcal B^c}D^s]}{\E V^s}
       &\le K\E_\nu[\one_{\mathcal B^c}d_{\min}^{-s}]\\
       &\le K(\E_\nu d_{\min}^{-2s})^{1/2}
                    \nu(\mathcal B^c)^{1/2}
        \le K R e^{-cR/2}.
\end{align*}
Both quantities tend to zero uniformly over all the prescriptions.
Combining these estimates with \Cref{boot:twist}, we obtain
\begin{equation}\label{boot:one-step}
 \E\norm{S^{\rm old}_{Y'X}}_{\rm HS}^s
 \le K_s(1+C_0)
       \left[e^{-scC_0}+(\E U^2)^{s/2}+\epsilon_R\right]\E V^s,
 \qquad \epsilon_R\longrightarrow0.
\end{equation}
The constants \(K_s,c>0\) are independent of large \(C_0\). First
choose \(C_0\) to make its exponential term small, then use
\Cref{boot:return} with \(n,m,N\) in order to make the return term
small, and enlarge \(m,N\) as necessary for \(\epsilon_R\). Thus the
entire multiplier can be made arbitrarily small. These are uniform
expectation bounds for each choice of mask and centers; no simultaneous
good event over all choices is needed.

\subsection{Contraction over all masks}

The passage from finite-scale smallness to exponential fractional-moment
decay parallels the finite-volume criteria of
\cite{AizenmanSchenkerFriedrichHundertmark2001}. We prove the contraction
here for the probe transmission moments and deterministic masks.

\begin{proposition}[Uniform mask contraction]\label{boot:contraction}
For fixed \(h,E,s\), parameters through \(R\) and an integer \(N_0\)
can be chosen so that, for every \(N\ge N_0\),
\[
 F_N(r)\le\tfrac12 F_N(r-4R),\qquad r\ge10R.
\]
Here \(F_N(r)\) is the supremum of
\(\E\norm{S_{Y'X}}_{\rm HS}^s\) over all deterministic artificial
masks, all subsets \(X,Y'\) lying in radius-\(R\) squares, and all
centers at distance at least \(r\). A supremum over an empty family
is zero. Consequently there are constants \(K,a_1>0\), independent
of \(N\ge N_0\), such that
\begin{equation}\label{boot:mask-decay}
                 F_N(r)\le K R^s e^{-a_1 r/R},\qquad r\ge0.
\end{equation}
\end{proposition}

\begin{proof}
Partition \(C\) into at most \(J\) sets \(C_k\), each contained in a
radius-\(R\) square with center in \(C\); \(J\) is an absolute
constant. Then
\[
 V^s=\left(\sum_{k=1}^J
            \norm{S^{\rm fill}_{C_kY'}}_{\rm HS}^2\right)^{s/2}
       \le\sum_{k=1}^J\norm{S^{\rm fill}_{C_kY'}}_{\rm HS}^s.
\]
Each new center is at distance at least \(r-4R\) from \(y\).
The filled model is included in the same mask class. Choose the
multiplier in \Cref{boot:one-step} at most \(1/(2J)\) and take
suprema, using symmetry to interchange the two block roles when
necessary. This proves the contraction. Always
\(F_N(r)\le K R^s\) by unitarity and the rank of a radius-\(R\) input
set. Iterating in steps of \(4R\) until the argument falls below
\(10R\) proves \Cref{boot:mask-decay}. All thresholds were fixed before
taking the supremum and are valid on every sufficiently large torus.
\end{proof}

\subsection{Removing the two endpoint probes}

Choose a mask consisting only of two constant-strength probes
\(\tau\) at distinct sites \(x,y\). Let
\(G^0=(H_N-E)^{-1}\), and let \(g\) be its two-by-two compression to
\(\{x,y\}\). Resolvent subtraction gives the damped compression
\[
                    g^{\rm damp}=(1-i\tau g)^{-1}g.
\]
For a symmetric two-by-two matrix, direct multiplication by the
adjugate shows
\[
 (g^{\rm damp})_{xy}
       =\frac{g_{xy}}{\det(1-i\tau g)}.
\]
Since \(S_{xy}=2i\tau(g^{\rm damp})_{xy}\), the exact identity is
\begin{equation}\label{boot:endpoint-identity}
          G^0_{xy}=\frac{S_{xy}}{2i\tau}\det(1-i\tau g).
\end{equation}

For completeness, the moments needed to remove the determinant are
uniform in the exterior and the volume. Conditional on all other
potentials, Schur complementation gives
\[
                    g=(A+\diag(v_x,v_y))^{-1}
\]
with \(A\) a fixed real symmetric matrix. Put
\(t=(v_x+v_y)/2\), \(u=(v_x-v_y)/2\). The integration square becomes
\(\abs u\le h\), \(\abs t\le h-\abs u\), with Jacobian two.
For each fixed \(u\), let \(\lambda_1(u),\lambda_2(u)\) be the real
eigenvalues of \(A+\diag(u,-u)\). If \(0<q<1\),
\[
 \int_{-(h-\abs u)}^{h-\abs u}
       \norm{(A+\diag(u,-u)+t1)^{-1}}^q\,\dd t
 \le\sum_{j=1}^2\int_{-h}^h\abs{t+\lambda_j(u)}^{-q}\,\dd t
 \le K_q h^{1-q}.
\]
Integration in \(u\), with the normalized joint uniform density,
therefore proves
\begin{equation}\label{boot:two-site}
                    \E\norm g^q\le K_{h,q},\qquad 0<q<1.
\end{equation}
The bound is uniform over all real symmetric \(A\), so no exterior
norm enters. The deleted system is generically invertible; its
exceptional set has zero product measure. A one-site Schur complement
and the same scalar integral also give
\(\E\abs{G^0_{xx}}^q\le K_{h,q}\).

Since \(\abs{\det(1-i\tau g)}\le(1+\tau\norm g)^2\),
\Cref{boot:two-site} with \(q=2s<1\) gives a uniform bound on
\(\E\abs{\det(1-i\tau g)}^s\). Apply Cauchy--Schwarz to
\Cref{boot:endpoint-identity}:
\[
 \E\abs{G^0_{xy}}^{s/2}
 \le(2\tau)^{-s/2}
       \bigl(\E\abs{S_{xy}}^s\bigr)^{1/2}
       \bigl(\E\abs{\det(1-i\tau g)}^s\bigr)^{1/2}.
\]
The first factor of expectations decays exponentially by
\Cref{boot:mask-decay}, applied to singleton subsets. The diagonal is
covered by the one-site bound. The now fixed \(R,\tau\) and all
constants are absorbed into \(A(h,E,s)\) and \(a(h,E,s)>0\), proving
\Cref{boot:fractional}. This conclusion remains a fixed-real-energy
finite-volume estimate; the spectral passage uses it separately at
each energy.
